% ============================================================
% ANHANG C – Literaturverzeichnis
% ============================================================

\section{Anhang~C: Literaturverzeichnis}
\label{sec:literaturverzeichnis}

\begin{hinweis}[Lehrhinweis]
  Dieses Literaturverzeichnis enthält alle Quellen, auf die im Dokument explizit
  verwiesen wird, sowie empfohlene Standardwerke für die im Modul behandelten
  Themen.
\end{hinweis}

% ── Architektur & arc42 ─────────────────────────────────────
\subsection*{Architektur \& arc42}

\begin{description}[leftmargin=2em, style=nextline]
  \item[Starke, G. (2024)]
  \textit{Effektive Softwarearchitekturen -- Ein praktischer Leitfaden.}\\
  Carl Hanser Verlag, 10.~Aufl., 2024. 369~S.,
  ISBN~978-3-446-47672-1.\\
  \textbf{Relevant: gesamtes Dokument.}

  \item[arc42.org (2025)]
  \textit{arc42 Template v9 (Juli 2025).}\\
  \url{https://arc42.org} -- Offizielle Template-Quelle.

  \item[Brown, S. (2018)]
  \textit{The C4 Model for Software Architecture.}\\
  \url{https://c4model.com} -- Grundlage der Diagrammnotation in Kap.~3, 5, 7.
\end{description}

% ── Domain-Driven Design ────────────────────────────────────
\subsection*{Domain-Driven Design}

\begin{description}[leftmargin=2em, style=nextline]
  \item[Evans, E. (2003)]
  \textit{Domain-Driven Design: Tackling Complexity in the Heart of Software.}\\
  Addison-Wesley. -- Grundlage für Bounded Contexts, Aggregates, Ubiquitous Language.
  \textbf{Relevant: Kap.~4, ADR-001.}

  \item[Vernon, V. (2013)]
  \textit{Implementing Domain-Driven Design.}\\
  Addison-Wesley. -- Praxishandbuch zu DDD-Mustern inkl.\ Context Maps und Aggregates.
  \textbf{Relevant: Kap.~4.2, 5.}

  \item[Brandolini, A. (2021)]
  \textit{Introducing EventStorming.}\\
  Leanpub. \url{https://www.eventstorming.com} --
  Methode zur kollaborativen Domänenentdeckung. \textbf{Relevant: Kap.~4.2, Anhang~A.}
\end{description}

% ── Microservices ───────────────────────────────────────────
\subsection*{Microservices}

\begin{description}[leftmargin=2em, style=nextline]
  \item[Newman, S. (2021)]
  \textit{Building Microservices: Designing Fine-Grained Systems.} 2.~Aufl.\\
  O'Reilly. -- Standardwerk zu Microservice-Architektur und Service-Kommunikation.
  \textbf{Relevant: Kap.~4.1 (Fussnote), ADR-001, ADR-002.}

  \item[Richardson, C. (2018)]
  \textit{Microservices Patterns: With Examples in Java.}\\
  Manning. \url{https://microservices.io} --
  Transactional Outbox Pattern, Saga, CQRS.
  \textbf{Relevant: ADR-004, ADR-009, Kap.~5.4.}
\end{description}

% ── Software-Qualität ───────────────────────────────────────
\subsection*{Software-Qualität}

\begin{description}[leftmargin=2em, style=nextline]
  \item[ISO/IEC 25010:2011]
  \textit{Systems and software engineering -- Systems and software Quality Requirements
  and Evaluation (SQuaRE) -- System and software quality models.}\\
  International Organization for Standardization. --
  Qualitätsmodell mit 8 Hauptmerkmalen.
  \textbf{Relevant: Kap.~1.2, 4.4, 10.}

  \item[Bass, L. / Clements, P. / Kazman, R. (2021)]
  \textit{Software Architecture in Practice.} 4.~Aufl.\\
  Addison-Wesley. -- Qualitätsszenarien und ATAM-Methode.
  \textbf{Relevant: Kap.~10, ATAM-light.}
\end{description}

% ── Sicherheit & Protokolle ─────────────────────────────────
\subsection*{Sicherheit \& Protokolle}

\begin{description}[leftmargin=2em, style=nextline]
  \item[IETF RFC 6749 (2012)]
  \textit{The OAuth 2.0 Authorization Framework.}\\
  \url{https://datatracker.ietf.org/doc/html/rfc6749} --
  Grundlage des Autorisierungskonzepts. \textbf{Relevant: ADR-003, Kap.~8.1.}

  \item[OpenID Foundation (2014)]
  \textit{OpenID Connect Core 1.0.}\\
  \url{https://openid.net/specs/openid-connect-core-1_0.html} --
  OIDC-Protokoll für Authentifizierung via externem IdP.
  \textbf{Relevant: ADR-003, Kap.~8.1.}

  \item[OWASP Foundation (2021)]
  \textit{OWASP Top 10 -- 2021.}\\
  \url{https://owasp.org/Top10/} --
  Sicherheits-Checkliste für den Security-Audit vor Go-Live.
  \textbf{Relevant: Kap.~1.1 (Go-Live-Kriterien v1.0), Q5, Q7.}
\end{description}

% ── Patterns & Tools ────────────────────────────────────────
\subsection*{Patterns \& Tools}

\begin{description}[leftmargin=2em, style=nextline]
  \item[Nygard, M. (2011)]
  \textit{Documenting Architecture Decisions.}\\
  \url{https://cognitect.com/blog/2011/11/15/documenting-architecture-decisions} --
  Ursprüngliches ADR-Format. \textbf{Relevant: Kap.~2 (Fussnote), Kap.~9.}

  \item[Fowler, M. (2018)]
  \textit{Strangler Fig Application.}\\
  \url{https://martinfowler.com/bliki/StranglerFigApplication.html} --
  Migrationsmuster; relevant für v2.0 Multi-Tenancy (Kap.~14.3).

  \item[Fowler, M. / Foemmel, M. (2006)]
  \textit{Continuous Integration.}\\
  \url{https://martinfowler.com/articles/continuousIntegration.html} --
  Grundlage der CI/CD-Strategie (Kap.~7.3).
\end{description}